% !TEX root =  mosaic.tex
\section{Introduction}
\label{sec:intro}

\gpus are the dominant execution platform for scientific computing, machine learning, and data analytics. To target them, modern compilers automatically apply loop transformations such as fusion, distribution, interchange, skewing, and tiling: polyhedral frameworks~\cite{ppcg,pluto,polly,tc,akg} derive these transformations through dependence analysis, while deep learning compilers like Halide~\cite{halide} and TVM~\cite{tvm} expose scheduling primitives that can be explored or tuned.

Transformation scheduling has been extensively explored in these systems, but efficient \gpu code generation also depends on architecture-specific execution parameters that remain after, or are chosen alongside, scheduling. Once a schedule and its mapping structure are fixed, the compiler must still choose tile sizes and thread-block dimensions that govern how the transformed computation consumes shared memory, registers, caches, and streaming multiprocessors (SMs). Unlike loop transformations, which reshape the iteration space and computation order, execution-space optimization leaves them intact and selects resource parameters exposed by the resulting schedule. These decisions therefore complement transformation scheduling rather than replace it.

Prior work on these execution parameters follows two broad lines: analytical selection and empirical search. Analytical tile size selection originated on CPUs, choosing tiles so a footprint stays resident in cache~\cite{mckinley-tss,lam-tss} and later embedding this reasoning in polyhedral~\cite{pluto-tss}, geometric programming~\cite{rajopadhye-gp,rajopadhye-igp}, and decoupled~\cite{tilemind} formulations. \gpu-specific analytical approaches include models derived for individual kernel classes such as matrix multiplication, convolution, or BLAS~\cite{li-gemm,saday-conv,blis}, compiler-integrated tile selection~\cite{ppcg-tss}, and EATSS~\cite{eatss}, which analytically selects \gpu tile sizes under a joint latency--energy objective. Empirical approaches, including autotuners~\cite{metaschedule,ansor,automatichalide} and learning-based pruning~\cite{hctuner,shirako-tss,yuki-tss}, instead explore candidate configurations through measurement or prediction and incur repeated tuning costs for new targets or shapes.

We formulate execution-space selection on a fixed transformed schedule: choosing tile sizes and thread-block dimensions while reasoning over abstract shared-memory placement scenarios and coupled \gpu resource effects. \name uses a compiler-agnostic metadata interface populated by compiler-specific adapters. The execution-space unknowns are coupled: the per-thread register budget depends on the resident-block count, which itself depends on register usage; staging an operand in shared memory cuts global traffic but consumes capacity and can lower the resident-block count; an SM hosts an integer number of blocks, so the grid runs in a discrete number of waves whose boundaries move with tile and placement choices; and the exact register count is fixed by the backend allocator only after code generation.

Our key insight is that these interactions can be captured by two compact quantities whose interaction is what makes the landscape irregular: a \emph{quantized wave cost}, a per-block cost multiplied by an integer wave count, which reproduces the step-function discontinuities that smooth objectives miss, and a \emph{shared-read-per-FMA ratio}, which normalizes operand reads by the arithmetic they serve and exposes register blocking as the lever trading locality against register pressure. Around these, a conservative register-working-set filter rejects only configurations whose register lower bound already exceeds the budget, breaking the occupancy--register cycle by reasoning rather than measurement.

We realize this insight in \name, which operates on the transformed programs emitted by existing schedulers~\cite{ppcg,tvm}. \name selects tile sizes and thread-block dimensions while reasoning over abstract shared-memory placement scenarios and accounting for occupancy and register pressure. It enumerates these scenarios around a two-stage analytical selection with a quantized primary objective and a relaxed-cost tie-break over the Stage-1 optimal discrete configurations, recovering balanced register blocking for outer-product contractions as an analytical outcome rather than an imposed rule. \name sits strictly between transformation scheduling and backend code generation: an upstream compiler first fixes a legal transformation schedule and mapping structure---i.e., which dimensions are assigned to block and thread axes---while \name selects the tile sizes and thread-block dimensions that this structure leaves open. It neither searches nor modifies the schedule or mapping structure. Our current cost model targets conventional CUDA-core SIMT execution, as exercised by the PPCG and non-tensorized TVM paths evaluated in this work; we return to this scope boundary in \autoref{sec:conclusion}.

This distinguishes \name from the prior approaches above: unlike heuristic polyhedral code generation~\cite{ppcg}, it selects execution parameters through an architecture-aware analytical model; unlike kernel-specific \gpu models~\cite{li-gemm,saday-conv,blis}, it spans reductions, contractions, and stencils; unlike prior analytical tile-size selection approaches~\cite{ppcg-tss,eatss,pluto-tss}, it combines placement-aware analytical reasoning with register feasibility and wave quantization; and unlike autotuners~\cite{ansor,metaschedule}, it derives interpretable configurations once, with no per-shape or per-target tuning loop. The evaluation in \autoref{sec:eval} confirms that this analytical reasoning replaces a substantial portion of \gpu execution-space autotuning.

In summary, this work makes the following contributions:
\begin{itemize}
  \item We formulate \gpu execution-space selection as a distinct post-scheduling optimization problem in GPU code generation, complementary to transformation scheduling.
  \item We propose \name, a unified analytical model that selects tile sizes and thread-block dimensions while reasoning over abstract shared-memory placement scenarios and accounting for occupancy and register pressure, closing the coupled feedback loop in one solvable formulation and yielding balanced register blocking as a closed-form optimum for outer-product contractions.
  \item We design a compiler-agnostic analytical core with compiler-specific metadata adapters that operates on transformed programs and integrates with polyhedral schedulers, template-based code generators, and autotuning systems.
  \item We evaluate \name with a cross-generation PolyBench/C~\cite{polybench} study on three NVIDIA \gpus, together with deep learning operators and end-to-end workloads on the RTX 4090, showing aggregate gains over heuristics and analytical selectors and competitive performance with a state-of-the-art autotuner without measurement-driven execution-space search.
\end{itemize}
